Je werkt als junior systeembeheerder in een fictief bedrijf. De ICT-manager wil dat je aantoont dat je veilig kunt werken met herstelpunten in Windows 11, voordat je updates of nieuwe software installeert.

Je krijgt de opdracht om:
\begin{enumerate}
\item Theorie - Leg in eigen woorden uit wat een herstelpunt doet en wat de verschillen zijn met:
  \begin{itemize}
  \item een back-up
  \item een systeemimage
  \end{itemize}
\item Praktijk
  \begin{enumerate}
  \item Controleer of systeembeveiliging aanstaat.
  \item Maak een nieuw herstelpunt met een duidelijke naam.
  \item Documenteer met schermafbeeldingen hoe je dit hebt gedaan.
  \end{enumerate}
\item Testen
  \begin{enumerate}
  \item Installeer een klein programma (of wijzig een instelling.)
  \item Zet het systeem terug met behulp van het gemaakte herstelpunt.
  \item Controleer of de wijziging ongedaan is gemaakt.
  \end{enumerate}
\item Reflectie \& Rapportage - Maak een kort verslag (max. 1 A4) met:
  \begin{itemize}
  \item Stap-voor-stap uitleg (incl. Screenshots) van wat je gedaan hebt.
  \item Resultaten van de test.
  \item Voor- en nadelen van herstelpunten in de praktijk.
  \end{itemize}
\end{enumerate}

